\documentclass[11pt,reqno]{amsart}

\usepackage[a4paper,margin=1.15in]{geometry}
\usepackage[T1]{fontenc}
\usepackage{lmodern}
\usepackage{amsmath,amssymb,mathtools}
\usepackage[hidelinks]{hyperref}

\newtheorem{question}{Question}
\newtheorem{assumption}{Assumption}
\newcommand{\cH}{\mathcal H}
\newcommand{\cM}{\mathcal M}
\newcommand{\cZ}{\mathcal Z}
\newcommand{\id}{\mathbf 1}

\title{Question I001: Full Tracial Corners and Finite Central Fibers}
\date{August 2, 2026}

\begin{document}
\maketitle

\begin{question}
Let $\cM\subseteq B(\cH)$ be a concrete von Neumann algebra and let
$q\in\cM$ be a projection.  Suppose that $q\cM q$ carries a faithful
tracial state and that $q$ is full, in the explicit sense that
\begin{equation}
  0\neq z\in\cZ(\cM)\quad\Longrightarrow\quad zq\neq0.
  \label{eq:question_fullness}
\end{equation}
Can one formalize, using mathlib's concrete
\texttt{VonNeumannAlgebra H}, the conclusion that every nonzero central
summand $z\cM$ contains a nonzero finite projection?
\end{question}

\begin{assumption}[Finite-fiber specialization]
The center is presented by finitely many nonzero pairwise orthogonal central
projections $(z_i)_{i\in I}$ with $\sum_{i\in I}z_i=\id$.  When each
$z_i\cM$ is a factor, it is called a finite central fiber.
\end{assumption}

The immediate Lean target is the stronger ambient central-summand statement:
it supplies a nonzero finite projection $e\leq z_i$ for every index $i$.
The standard factor/type-III conclusion requires a later bundled construction
of the represented central corner $z_i\cM$ and is deliberately not claimed as
machine checked here.  The central-summand statement is primary because it is
precisely the conclusion of the current manuscript's
Theorem~\texttt{thm:finite\_projections\_exist}.

\end{document}
